% The Long Tail Is Unaudited — arXiv submission source.
% Compile: paste into Overleaf (pdfLaTeX, article class). Numbers: live
% Supabase snapshot 2026-07-08. Refresh via the queries in
% ops/research/intake-launch-kit.md if submission slips by weeks.
\documentclass[11pt]{article}
\usepackage[margin=1in]{geometry}
\usepackage[hidelinks]{hyperref}
\usepackage{booktabs}
\usepackage{microtype}

\title{The Long Tail Is Unaudited:\\A Security Census of 131,933 Open-Source AI Agent Skills and MCP Servers}
\author{Jason Zhu\\Independent researcher \quad \href{https://agentskillshub.top}{agentskillshub.top} \quad \texttt{m17551076169@gmail.com}}
\date{Snapshot: July 8, 2026}

\begin{document}
\maketitle

\begin{abstract}
AI coding agents (Claude Code, Cursor, OpenAI Codex, and others) increasingly
install third-party ``skills'' and Model Context Protocol (MCP) servers that
execute with the agent's full permissions and often its credentials. Yet the
marketplaces that distribute these artifacts rank them by popularity and
recency, not safety, and expose no per-item trust signal. We collect
\textbf{131,933} open-source agent skills and MCP servers from GitHub and grade
them with a rule-based scanner spanning 35 red-flag patterns derived from an
established agent-security taxonomy. Of the \textbf{20,368} artifacts popular
enough to grade ($\geq$5 stars), \textbf{3.3\%} are graded UNSAFE or REJECT
($\approx$1 in 31) and \textbf{8.9\%} carry at least a CAUTION flag. Our
central finding is distributional: the UNSAFE rate is \textbf{4.5\% in the
low-star long tail (5--20 stars) versus 0.4\% among popular repositories
(1,000+ stars)} --- an $\sim$11$\times$ gap --- while \textbf{85\% of the full
catalog is unaudited}. Risk concentrates exactly where popularity ranking
hides it and where no marketplace provides a signal. We release the complete
graded dataset (131,933 rows, CC~BY~4.0) and the scanner rules.
\end{abstract}

\section{Introduction}

An agent skill is a package of instructions, and often code, that teaches an
AI agent to perform a task; an MCP server is a process that exposes tools to
an agent over a standard protocol. Both run inside the trust boundary of the
user's agent: a skill's install script runs on the user's machine, and an MCP
server receives whatever credentials the agent holds. Installing one is a
supply-chain decision, but it is presented to users as a one-click
convenience.

The marketplaces that have emerged (one indexes over 300K skills; others range
from thousands to millions of monthly visits) rank by stars, installs, or
recency. None, to our knowledge, ships a per-item security grade. That the
single most-installed artifact on one such marketplace is a community-built
``skill-vetter'' is direct evidence that users feel this gap and are
improvising around it.

We ask three empirical questions. (Q1)~How prevalent are unsafe artifacts
across the open-source agent-skill ecosystem? (Q2)~How is that risk
distributed with respect to popularity --- does the community's star signal
already protect users? (Q3)~What categories of risk dominate? Our
contributions:

\begin{enumerate}
  \item The broadest census of the published agent-skill ecosystem to date:
        131,933 artifacts, refreshed on an 8-hour cycle.
  \item An open, rule-based grader over 35 red-flag patterns in three
        severity tiers (Appendix~A).
  \item A released, citable dataset of grades, red-flags, and quality scores
        --- to our knowledge the first full graded dataset at this scale.
\end{enumerate}

\section{Background, Threat Model, and Related Work}

\textbf{Artifacts.} We consider skills (Claude Skills, Codex skills, and
portable \texttt{SKILL.md} variants), MCP servers, and standalone agent
tools. All can carry executable code (install hooks, Python entry points) and
all are loaded into a running agent's context.

\textbf{Threat classes.} We group risks into six families, mapped to concrete
red-flags: credential exfiltration (reading \texttt{.env}, agent
config/memory, tokens), obfuscated or remote code execution
(\texttt{curl|sh}, \texttt{eval}, base64-decoded exec), persistence (cron,
service installation, shell-rc injection), tunneling/egress (reverse shells,
tunnel services), privilege (unexpected \texttt{sudo}, privileged
containers), and supply-chain (runtime install-and-execute, typosquatting
analogues).

\textbf{Related work.} Agent-skill and MCP security has drawn rapid attention
in 2026. Closest to ours, Liu et al.\ (arXiv:2601.10338) evaluate 31,132
skills (of 42,447 collected) with the SkillScan framework and report a 26.1\%
vulnerability rate; Etteib, Lunghi, and Bissyand\'e (arXiv:2606.23416) build
an attention-based detector that locates malicious instruction spans in
marketplace skills. A second line targets enforcement and formal guarantees:
SkillScope (arXiv:2605.05868) proposes fine-grained least-privilege
enforcement, while a formal supply-chain analysis (arXiv:2603.00195) and
audit--runtime-gap work (arXiv:2605.05274) model the trust boundary
explicitly. A third studies the MCP protocol itself: threat modeling and
tool-poisoning (arXiv:2603.22489) and a defense-placement taxonomy
(arXiv:2604.07551). These contribute detection methods, enforcement
mechanisms, and threat taxonomies. Our contribution is orthogonal and
complementary: the broadest static census of the \emph{published} ecosystem,
a \emph{released} graded dataset, and a distributional finding --- unsafe
density is an order of magnitude higher in the unranked long tail. We further
relate to npm/PyPI typosquatting and compromised-dependency studies as
supply-chain analogues.

\section{Method}

\textbf{Collection.} A six-phase pipeline runs every 8 hours: (1)~GitHub
search across a fixed set of query patterns for the artifact types;
(2)~verified-creator repository fetch; (3)~community submissions;
(4)~owner/metadata enrichment; (5)~README fetch ($\leq$50\,KB); (6)~dedup by
\texttt{owner/repo} and upsert. This yields 131,933 distinct artifacts at the
July 8, 2026 snapshot.

\textbf{Grading.} A static scanner matches each artifact's README, install
scripts, configuration, and file tree against 35 red-flag patterns in three
severity tiers (Appendix~A), producing a grade: SAFE (no risky pattern),
CAUTION (accumulated medium-tier flags), UNSAFE (a high-tier pattern), or
REJECT (a confirmed-malicious pattern). Artifacts below a popularity gate
($<$5 stars) or lacking scannable content are marked UNAUDITED. We also
compute a 0--100 quality score from documentation, structure, and maintenance
signals (orthogonal to security; included in the dataset for context).

\textbf{Honesty of the instrument.} The scanner is static and point-in-time.
It is a \emph{floor, not a ceiling}: it catches declared patterns, not
obfuscated intent or runtime behavior, and does not execute code. It has not
yet been validated against a hand-labeled ground truth, so we present the
percentages below as \emph{descriptive statistics of the released dataset},
not as precise prevalence estimates. Two considerations bound the impact of
this limitation. First, our central finding is \emph{relative}: unless the
scanner's error rate varies systematically with repository popularity ---
and we know of no mechanism by which regex matching would --- the
$\sim$11$\times$ long-tail gap stands even if the absolute level shifts.
Second, our numbers are not comparable to runtime evaluations: Liu et al.'s
26.1\% measures vulnerability under adversarial prompting at runtime, while
our 3.3\% measures statically declared dangerous patterns; the two measure
different phenomena and can both be true. UNAUDITED means \emph{not
assessed} --- it must not be read as ``safe.'' A validated, layered audit
(LLM-assisted review plus sandboxed execution) is future work; the released
dataset is the substrate for it.

\section{Results}

All numbers are from the live July 8, 2026 snapshot.

\textbf{Prevalence (Q1).} Of 131,933 indexed artifacts, 20,368 are graded
(15.4\%). Among graded: SAFE 18,563; CAUTION 1,141; UNSAFE 639; REJECT 25.
Thus 664 (3.3\%, $\approx$1 in 31) are UNSAFE or REJECT and 1,805 (8.9\%) are
CAUTION-or-worse.

\textbf{Distribution by popularity (Q2) --- the central result.} The
UNSAFE/REJECT rate falls monotonically with popularity
(Table~\ref{tab:buckets}).

\begin{table}[h]
\centering
\begin{tabular}{lrrr}
\toprule
Star bucket & $n$ & UNSAFE/REJECT & rate \\
\midrule
5--20      & 8,386 & 374 & \textbf{4.5\%} \\
20--100    & 6,160 & 243 & 3.9\% \\
100--1,000 & 4,370 & 41  & 0.9\% \\
1,000+     & 1,452 & 6   & \textbf{0.4\%} \\
\bottomrule
\end{tabular}
\caption{UNSAFE/REJECT rate by popularity. The low-star tail carries
$\sim$11$\times$ the unsafe density of the popular head.}
\label{tab:buckets}
\end{table}

Popularity is weakly protective --- but the corollary is that \textbf{85\% of
the catalog (111,565 artifacts) is UNAUDITED}, and the residual unsafe
density is $\sim$11$\times$ higher in the low-star tail that no marketplace
ranks or flags. The risk is precisely where users get the least signal.

\textbf{Risk categories (Q3).} Among flagged artifacts, the most common
red-flags are unexpected privilege (\texttt{sudo\_usage}, 1,114), service
persistence (309), \texttt{curl|sh} remote execution (228), agent-config
theft (148), tunneling (132), sensitive environment-variable access (103),
and \texttt{eval} (99). Credential- and persistence-oriented flags --- the
ones that matter most when an artifact runs with agent credentials --- are
well represented.

\textbf{By category.} UNSAFE rates are comparable across artifact types: MCP
servers 3.5\% ($n$=8,775), Claude skills 3.9\% ($n$=3,320), Codex skills
3.4\% ($n$=2,614), agent tools 2.6\% ($n$=4,659). No single artifact class is
disproportionately dangerous; the problem is ecosystem-wide.

\section{Discussion}

The headline is not ``3.3\% are unsafe'' but \emph{where} the 3.3\% lives.
Popularity ranking, the default on every marketplace, systematically routes
users toward the safer head and provides no warning in the long tail where
density is an order of magnitude higher. A per-install trust signal --- a
grade shown before the install command --- directly addresses this, and costs
the user nothing to read.

Static grading has clear limits (\S3). It motivates layered auditing:
LLM-assisted source review for intent, and sandboxed execution for behavior.
We frame the released dataset as the substrate for that work: it provides the
population, the sampling frame, and a baseline signal against which deeper
audits can be measured.

\section{Availability}

\begin{itemize}
  \item \textbf{Dataset:}
        \href{https://huggingface.co/datasets/jasonzhuyansen/agent-skills-security-grades}{huggingface.co/datasets/jasonzhuyansen/agent-skills-security-grades}
        (131,933 rows; parquet + CSV; CC~BY~4.0).
  \item \textbf{Archived snapshot (DOI):}
        \href{https://doi.org/10.5281/zenodo.21292799}{doi:10.5281/zenodo.21292799}
        (Zenodo, CSV snapshot of 2026-07-06).
  \item \textbf{Living version:} \href{https://agentskillshub.top}{agentskillshub.top},
        regraded every 8 hours, with per-repository audit pages.
  \item \textbf{Scanner rules:} the full pattern set is enumerated in
        Appendix~A; source is open on GitHub.
\end{itemize}

\appendix
\section{Red-Flag Patterns}

The scanner evaluates 35 patterns in three tiers. A REJECT-tier match forces
REJECT; a HIGH-tier match forces UNSAFE; MEDIUM-tier matches accumulate to
CAUTION (2+). Verified-organization allowlists (anthropics, openai, google,
\ldots) raise vendor trust and are not red-flags.

\subsection*{REJECT tier (confirmed malicious, 2)}
\small
\begin{tabular}{lp{9cm}}
\texttt{exfil\_secrets\_combo} & exfiltrates secrets via pipe to a network tool \\
\texttt{backdoor\_install} & installs backdoor via shell startup + remote download \\
\end{tabular}

\subsection*{HIGH tier ($\rightarrow$ UNSAFE, 23)}
\small
\begin{tabular}{lp{9cm}}
\texttt{data\_exfiltration} & sends local data to external server via curl POST \\
\texttt{curl\_pipe\_shell} & downloads and executes remote script via \texttt{curl|bash} \\
\texttt{wget\_pipe\_shell} & downloads and executes remote script via \texttt{wget|bash} \\
\texttt{credential\_harvest} & harvests credentials from environment variables \\
\texttt{env\_file\_read} & reads \texttt{.env} files which may contain secrets \\
\texttt{sensitive\_dir\_access} & accesses \texttt{\textasciitilde/.ssh}, \texttt{\textasciitilde/.aws}, \texttt{\textasciitilde/.gnupg} \\
\texttt{etc\_sensitive\_read} & reads \texttt{/etc/shadow}, \texttt{/etc/passwd} \\
\texttt{agent\_memory\_theft} & accesses agent memory/identity files \\
\texttt{agent\_config\_theft} & accesses agent configuration/session files \\
\texttt{exec\_import} & executes dynamically imported Python code \\
\texttt{base64\_exec} & decodes and pipes base64 data for execution \\
\texttt{chmod\_dangerous} & sets dangerous permissions (777 or setuid) \\
\texttt{write\_etc} & writes to system \texttt{/etc/} \\
\texttt{privilege\_escalation} & attempts privilege escalation to root \\
\texttt{cron\_persistence} & installs cron job for persistence \\
\texttt{shell\_rc\_inject} & injects commands into shell startup files \\
\texttt{service\_persistence} & installs persistent service/daemon \\
\texttt{reverse\_shell} & opens reverse shell connection \\
\texttt{dev\_tcp} & uses \texttt{/dev/tcp} (reverse-shell indicator) \\
\texttt{rm\_rf\_root} & recursively deletes from root filesystem \\
\texttt{obfuscated\_exec} & executes obfuscated/encoded Python code \\
\texttt{hex\_encoded\_payload} & hex-encoded payload (possible obfuscation) \\
\texttt{runtime\_install\_exec} & installs and immediately executes package at runtime \\
\end{tabular}

\subsection*{MEDIUM tier (2+ $\rightarrow$ CAUTION, 10)}
\small
\begin{tabular}{lp{9cm}}
\texttt{sudo\_usage} & uses sudo for elevated privileges \\
\texttt{docker\_privileged} & runs Docker container in privileged mode \\
\texttt{fs\_root\_access} & reads root filesystem directory \\
\texttt{sensitive\_env\_vars} & references multiple sensitive API keys/tokens \\
\texttt{ssl\_disabled} & disables SSL/TLS certificate verification \\
\texttt{eval\_usage} & uses \texttt{eval()} for dynamic code execution \\
\texttt{raw\_ip\_request} & HTTP request to raw IP address \\
\texttt{env\_access} & accesses environment variables programmatically \\
\texttt{subprocess\_spawn} & spawns subprocesses for command execution \\
\texttt{tunnel\_service} & uses tunneling service to expose local network \\
\end{tabular}

\section{Reproducibility}

Collection query patterns, the 5-star grading gate, scanner version, and the
snapshot date (2026-07-08) fully determine the released dataset. The dataset
schema is documented on the HuggingFace card; the site regrades on an 8-hour
cycle, so the living numbers drift from this snapshot by design.

\end{document}
